\documentclass[12pt,a4paper]{article}

% ── Packages ────────────────────────────────────────────────────────────────
\usepackage[utf8]{inputenc}
\usepackage[T1]{fontenc}
\usepackage{lmodern}
\usepackage[margin=2.5cm, headheight=15pt]{geometry}
\usepackage{amsmath, amssymb, amsthm}
\usepackage{mathtools}
\usepackage{enumitem}
\usepackage{xcolor}
\usepackage{titlesec}
\usepackage{fancyhdr}
\usepackage{hyperref}
\usepackage{tcolorbox}
\tcbuselibrary{skins,breakable}

% ── Colours ─────────────────────────────────────────────────────────────────
\definecolor{primary}{RGB}{0,70,127}
\definecolor{secondary}{RGB}{180,20,20}
\definecolor{solutionbg}{RGB}{244,250,244}
\definecolor{solutionborder}{RGB}{0,110,0}
\definecolor{exbg}{RGB}{240,245,255}

\hypersetup{colorlinks=true, linkcolor=primary}

% ── Section formatting ───────────────────────────────────────────────────────
\titleformat{\section}[block]
  {\large\bfseries\color{primary}}
  {\thesection.}{0.5em}{}[\titlerule]

% ── Page style ───────────────────────────────────────────────────────────────
\pagestyle{fancy}
\fancyhf{}
\lhead{\textcolor{primary}{\textbf{MTH253 -- Calculus III}}}
\rhead{\textcolor{primary}{Full Solutions}}
\cfoot{\thepage}
\renewcommand{\headrulewidth}{0.4pt}

% ── tcolorbox environments ───────────────────────────────────────────────────
\newtcolorbox{exbox}[1][]{%
  breakable, enhanced,
  colback=exbg, colframe=primary,
  fonttitle=\bfseries\color{primary},
  title={Exercise}, #1
}
\newtcolorbox{solbox}[1][]{%
  breakable, enhanced,
  colback=solutionbg, colframe=solutionborder,
  fonttitle=\bfseries\color{solutionborder},
  title={Solution}, #1
}

% ── Theorem style for exercise counter ───────────────────────────────────────
\newtheoremstyle{exstyle}
  {4pt}{4pt}{\normalfont}{0pt}{\bfseries\color{primary}}{.}{0.5em}{}
\theoremstyle{exstyle}
\newtheorem{exercise}{Exercise}[section]

% ── Macros ───────────────────────────────────────────────────────────────────
\newcommand{\R}{\mathbb{R}}
\newcommand{\dd}{\,\mathrm{d}}
\newcommand{\pd}[2]{\dfrac{\partial #1}{\partial #2}}
\newcommand{\grad}{\nabla}
\newcommand{\dif}[1]{\,\mathrm{d}#1}

% ────────────────────────────────────────────────────────────────────────────
\begin{document}

\begin{center}
  {\LARGE\bfseries\color{primary} Calculus III --- Full Solutions}\\[4pt]
  {\large\color{secondary}
   Chapter 1: Differentiation \quad|\quad Chapter 2: Multiple Integrals}\\[4pt]
  {\normalsize University of Science, VNU-HCM
   \hfill Academic Year 2024--2025}
\end{center}
\vspace{-4pt}\hrule height 1.2pt\vspace{10pt}

\tableofcontents
\newpage

% ============================================================
\section{Vectors and Geometry of Space}
% ============================================================

%------------------------------------------------------------------
\begin{exercise}
Find the vector $\overrightarrow{AB}$ and its magnitude for (a) $A(1,-2,4)$, $B(3,0,-1)$;
(b) $A(-3,5,0)$, $B(2,-1,6)$.
\end{exercise}
\begin{solbox}
\textbf{(a)} $\overrightarrow{AB}=\langle 3-1,\;0-(-2),\;-1-4\rangle=\langle 2,2,-5\rangle$.
\[
  |\overrightarrow{AB}|=\sqrt{4+4+25}=\sqrt{33}.
\]

\textbf{(b)} $\overrightarrow{AB}=\langle 2-(-3),\;-1-5,\;6-0\rangle=\langle 5,-6,6\rangle$.
\[
  |\overrightarrow{AB}|=\sqrt{25+36+36}=\sqrt{97}.
\]
\end{solbox}

%------------------------------------------------------------------
\begin{exercise}
Let $\mathbf{u}=\langle 2,-1,3\rangle$ and $\mathbf{v}=\langle -1,4,2\rangle$.
Find (a) $\mathbf{u}\cdot\mathbf{v}$; (b) angle between them; (c) $\mathbf{u}\times\mathbf{v}$;
(d) unit vector in direction of $\mathbf{u}$.
\end{exercise}
\begin{solbox}
\textbf{(a)}
$\mathbf{u}\cdot\mathbf{v}=2(-1)+(-1)(4)+3(2)=-2-4+6=\boxed{0}$.
The vectors are \emph{orthogonal}.

\textbf{(b)} Since $\mathbf{u}\cdot\mathbf{v}=0$, the angle is $\theta=\pi/2$.

\textbf{(c)}
\[
  \mathbf{u}\times\mathbf{v}
  =\begin{vmatrix}\mathbf{i}&\mathbf{j}&\mathbf{k}\\2&-1&3\\-1&4&2\end{vmatrix}
  =\mathbf{i}\bigl((-1)(2)-(3)(4)\bigr)-\mathbf{j}\bigl((2)(2)-(3)(-1)\bigr)+\mathbf{k}\bigl((2)(4)-(-1)(-1)\bigr)
\]
\[
  =\mathbf{i}(-2-12)-\mathbf{j}(4+3)+\mathbf{k}(8-1)=\langle -14,-7,7\rangle.
\]
\emph{Verification:} $\mathbf{u}\cdot\langle-14,-7,7\rangle=2(-14)+(-1)(-7)+3(7)=-28+7+21=0$\checkmark

\textbf{(d)} $|\mathbf{u}|=\sqrt{4+1+9}=\sqrt{14}$.
Unit vector $=\dfrac{\mathbf{u}}{|\mathbf{u}|}=\left\langle\dfrac{2}{\sqrt{14}},\dfrac{-1}{\sqrt{14}},\dfrac{3}{\sqrt{14}}\right\rangle$.
\end{solbox}

%------------------------------------------------------------------
\begin{exercise}
(a) Parametric and symmetric equations of the line through $P(-1,2,3)$ parallel to
$\mathbf{v}=\langle 3,-1,5\rangle$; (b) line through $A(1,0,-2)$ and $B(4,3,1)$;
(c) intersection of $x=1+2t,\;y=-t,\;z=3+t$ with the plane $2x-y+z=7$.
\end{exercise}
\begin{solbox}
\textbf{(a)}
\emph{Parametric:} $x=-1+3t,\;y=2-t,\;z=3+5t$.\quad
\emph{Symmetric:} $\dfrac{x+1}{3}=\dfrac{y-2}{-1}=\dfrac{z-3}{5}$.

\textbf{(b)} Direction $\mathbf{v}=\overrightarrow{AB}=\langle 3,3,3\rangle$.
Parametric: $x=1+3t,\;y=3t,\;z=-2+3t$.

\textbf{(c)} Substitute into the plane:
\[
  2(1+2t)-(-t)+(3+t)=7\implies 2+4t+t+3+t=7\implies 6t=2\implies t=\tfrac{1}{3}.
\]
Point: $x=\tfrac{5}{3}$, $y=-\tfrac{1}{3}$, $z=\tfrac{10}{3}$.
\end{solbox}

%------------------------------------------------------------------
\begin{exercise}
(a) Plane through $P(2,-1,3)$ with $\mathbf{n}=\langle 1,5,-2\rangle$;
(b) plane through $P(0,1,1)$, $Q(1,0,1)$, $R(1,1,0)$;
(c) Are $2x-3y+z=4$ and $4x-6y+2z=1$ parallel?
\end{exercise}
\begin{solbox}
\textbf{(a)} $(x-2)+5(y+1)-2(z-3)=0\implies x+5y-2z+9=0$.

\textbf{(b)} $\overrightarrow{PQ}=\langle 1,-1,0\rangle$, $\overrightarrow{PR}=\langle 1,0,-1\rangle$.
\[
  \mathbf{n}=\overrightarrow{PQ}\times\overrightarrow{PR}=
  \begin{vmatrix}\mathbf{i}&\mathbf{j}&\mathbf{k}\\1&-1&0\\1&0&-1\end{vmatrix}
  =\langle(-1)(-1)-0,\;0(1)-1(-1),\;0-(-1)\rangle=\langle 1,1,1\rangle.
\]
Plane: $(x-0)+(y-1)+(z-1)=0\implies x+y+z=2$.

\textbf{(c)} Rewrite the second plane: $2x-3y+z=\tfrac{1}{2}$.
Both have normal $\langle 2,-3,1\rangle$ but different constants ($4\neq\tfrac{1}{2}$), so they are \textbf{parallel} (and distinct).
\end{solbox}

%------------------------------------------------------------------
\begin{exercise}
Describe and sketch: (a) $x^2+z^2=9$; (b) $z=4-x^2-y^2$; (c) $\frac{x^2}{4}+\frac{y^2}{9}+\frac{z^2}{16}=1$.
\end{exercise}
\begin{solbox}
\textbf{(a)} Circular cylinder of radius $3$ with axis along the $y$-axis.

\textbf{(b)} Circular paraboloid opening downward, vertex at $(0,0,4)$.

\textbf{(c)} Ellipsoid with semi-axes $a=2$ (along $x$), $b=3$ (along $y$), $c=4$ (along $z$).
\end{solbox}

% ============================================================
\section{Functions of Several Variables}
% ============================================================

%------------------------------------------------------------------
\begin{exercise}
For each function find the domain, range, and level curves:
(a) $f(x,y)=\sqrt{16-x^2-y^2}$; (b) $f(x,y)=\ln(x+y-1)$;
(c) $f(x,y)=\dfrac{x^2-y^2}{x^2+y^2}$.
\end{exercise}
\begin{solbox}
\textbf{(a)} Domain: $x^2+y^2\le 16$ (disk of radius 4). Range: $[0,4]$.
Level curves $\sqrt{16-x^2-y^2}=k$: circles $x^2+y^2=16-k^2$, $0\le k\le 4$.

\textbf{(b)} Domain: $x+y>1$ (open half-plane). Range: $\R$.
Level curves: $x+y=1+e^k$ (parallel lines).

\textbf{(c)} Domain: $\R^2\setminus\{(0,0)\}$. Writing in polar: $f=\cos(2\theta)$, so Range $=[-1,1]$.
Level curves: $x^2(1-k)=y^2(1+k)$, i.e.\ $y=\pm x\sqrt{\tfrac{1-k}{1+k}}$ for $|k|<1$ (lines through origin); $k=1$ gives $y=0$; $k=-1$ gives $x=0$.
\end{solbox}

%------------------------------------------------------------------
\begin{exercise}
Find the domain and level surfaces of $f(x,y,z)=\sqrt{9-x^2-y^2-z^2}$.
\end{exercise}
\begin{solbox}
Domain: $x^2+y^2+z^2\le 9$ (closed ball of radius 3).
Level surfaces $f=k$: $x^2+y^2+z^2=9-k^2$, spheres of radius $\sqrt{9-k^2}$ for $0\le k\le 3$.
\end{solbox}

%------------------------------------------------------------------
\begin{exercise}
Sketch the level curves of $f(x,y)=x^2+4y^2$ for $k=0,1,4,9$.
\end{exercise}
\begin{solbox}
Setting $x^2+4y^2=k$:
\begin{itemize}[topsep=2pt,itemsep=1pt]
  \item $k=0$: single point $(0,0)$.
  \item $k=1$: ellipse $x^2/1+y^2/(1/4)=1$, semi-axes $a=1$, $b=\tfrac{1}{2}$.
  \item $k=4$: ellipse $x^2/4+y^2/1=1$, semi-axes $a=2$, $b=1$.
  \item $k=9$: ellipse $x^2/9+y^2/(9/4)=1$, semi-axes $a=3$, $b=\tfrac{3}{2}$.
\end{itemize}
The level curves are concentric ellipses expanding outward.
\end{solbox}

% ============================================================
\section{Limits and Continuity}
% ============================================================

%------------------------------------------------------------------
\begin{exercise}
Evaluate each limit or show it does not exist.
\end{exercise}
\begin{solbox}
\textbf{(a)} $\displaystyle\lim_{(x,y)\to(1,2)}\bigl(x^3y-2xy^2+5\bigr)$.
The function is a polynomial, hence continuous everywhere:
\[
  =1^3\cdot 2-2\cdot 1\cdot 4+5=2-8+5=-1.
\]

\textbf{(b)} $\displaystyle\lim_{(x,y)\to(0,0)}\frac{5x^2y}{x^2+y^2}$.
Using polar coordinates $x=r\cos\theta$, $y=r\sin\theta$:
\[
  \left|\frac{5x^2y}{x^2+y^2}\right|=\frac{5r^2\cos^2\theta\cdot|r\sin\theta|}{r^2}
  =5r|\cos^2\theta\sin\theta|\le 5r\to 0.
\]
By Squeeze Theorem, the limit is $\boxed{0}$.

\textbf{(c)} $\displaystyle\lim_{(x,y)\to(0,0)}\frac{x^2-xy}{\sqrt{x}-\sqrt{y}}$.
Factor: $x^2-xy=x(x-y)=x(\sqrt{x}-\sqrt{y})(\sqrt{x}+\sqrt{y})$. Hence
\[
  \frac{x^2-xy}{\sqrt{x}-\sqrt{y}}=x(\sqrt{x}+\sqrt{y})\to 0\cdot(0+0)=\boxed{0}.
\]

\textbf{(d)} $\displaystyle\lim_{(x,y)\to(0,0)}\frac{xy^2}{x^2+y^4}$.
Along $y=0$: limit $=0$. Along $x=y^2$: $\dfrac{y^2\cdot y^2}{y^4+y^4}=\dfrac{y^4}{2y^4}=\dfrac{1}{2}$.
Different limits $\Rightarrow$ limit \textbf{does not exist}.
\end{solbox}

%------------------------------------------------------------------
\begin{exercise}
Prove by the Squeeze Theorem that $\displaystyle\lim_{(x,y)\to(0,0)}\frac{x^2y^2}{x^2+y^2}=0$.
\end{exercise}
\begin{solbox}
Since $x^2\le x^2+y^2$, we have $x^2y^2\le(x^2+y^2)y^2$. Therefore
\[
  0\le\frac{x^2y^2}{x^2+y^2}\le y^2.
\]
Since $y^2\to 0$ as $(x,y)\to(0,0)$, the Squeeze Theorem gives the limit $=0$.\hfill$\square$
\end{solbox}

%------------------------------------------------------------------
\begin{exercise}
Discuss the continuity of
$f(x,y)=\begin{cases}2xy/(x^2+y^2)&(x,y)\neq(0,0)\\0&(x,y)=(0,0)\end{cases}$.
\end{exercise}
\begin{solbox}
For $(x,y)\neq(0,0)$: $f$ is a ratio of polynomials with nonzero denominator, hence continuous.

At $(0,0)$: along $y=0$, $\lim f=0$; along $y=x$, $\lim f=2x^2/(2x^2)=1\neq 0$.
The limit does not exist (or does not equal $f(0,0)=0$), so $f$ is \textbf{discontinuous at $(0,0)$}.
\end{solbox}

% ============================================================
\section{Partial Derivatives and Higher Derivatives}
% ============================================================

%------------------------------------------------------------------
\begin{exercise}
Find all first-order partial derivatives.
\end{exercise}
\begin{solbox}
\textbf{(a)} $f=x^4y-3x^2y^3+7y$:
\[
  f_x=4x^3y-6xy^3,\qquad f_y=x^4-9x^2y^2+7.
\]

\textbf{(b)} $g=e^x\sin(xy)$:
\begin{align*}
  g_x &=e^x\sin(xy)+e^x\cdot y\cos(xy)=e^x\bigl(\sin(xy)+y\cos(xy)\bigr),\\
  g_y &=e^x\cdot x\cos(xy)=xe^x\cos(xy).
\end{align*}

\textbf{(c)} $h=xy^2z^3+\ln(x+2y+3z)$:
\[
  h_x=y^2z^3+\frac{1}{x+2y+3z},\quad
  h_y=2xyz^3+\frac{2}{x+2y+3z},\quad
  h_z=3xy^2z^2+\frac{3}{x+2y+3z}.
\]

\textbf{(d)} $w=s\arctan(t/s)$. Let $u=t/s$; $\partial u/\partial s=-t/s^2$, $\partial u/\partial t=1/s$.
\begin{align*}
  w_s&=\arctan(t/s)+s\cdot\frac{1}{1+t^2/s^2}\cdot\frac{-t}{s^2}
       =\arctan(t/s)-\frac{t}{s^2+t^2},\\
  w_t&=s\cdot\frac{1}{1+t^2/s^2}\cdot\frac{1}{s}=\frac{s^2}{s^2+t^2}.
\end{align*}
\end{solbox}

%------------------------------------------------------------------
\begin{exercise}
Find all second-order partial derivatives of $f(x,y)=x^3y^2-2x^2y^4+5xy$
and verify Clairaut's theorem ($f_{xy}=f_{yx}$).
\end{exercise}
\begin{solbox}
First-order:
$f_x=3x^2y^2-4xy^4+5y$,\quad $f_y=2x^3y-8x^2y^3+5x$.

Second-order:
\begin{align*}
  f_{xx}&=6xy^2-4y^4,\\
  f_{yy}&=2x^3-24x^2y^2,\\
  f_{xy}&=\pd{}{y}(3x^2y^2-4xy^4+5y)=6x^2y-16xy^3+5,\\
  f_{yx}&=\pd{}{x}(2x^3y-8x^2y^3+5x)=6x^2y-16xy^3+5.
\end{align*}
Since $f_{xy}=f_{yx}=6x^2y-16xy^3+5$, Clairaut's theorem is verified.\hfill$\checkmark$
\end{solbox}

%------------------------------------------------------------------
\begin{exercise}
Verify that the following functions are harmonic ($f_{xx}+f_{yy}=0$).
\end{exercise}
\begin{solbox}
\textbf{(a)} $f=x^3-3xy^2$:
$f_{xx}=6x$, $f_{yy}=-6x$. Sum $=0$.\hfill$\checkmark$

\textbf{(b)} $f=e^x\cos y$:
$f_{xx}=e^x\cos y$, $f_{yy}=-e^x\cos y$. Sum $=0$.\hfill$\checkmark$

\textbf{(c)} $f=\ln(x^2+y^2)$:
\[
  f_x=\frac{2x}{x^2+y^2},\quad
  f_{xx}=\frac{2(x^2+y^2)-4x^2}{(x^2+y^2)^2}=\frac{2y^2-2x^2}{(x^2+y^2)^2}.
\]
By symmetry, $f_{yy}=\frac{2x^2-2y^2}{(x^2+y^2)^2}$.
$f_{xx}+f_{yy}=0$.\hfill$\checkmark$
\end{solbox}

%------------------------------------------------------------------
\begin{exercise}
Let $f(x,y)=x^2e^y+y\sin x$. Calculate $f_{xxy}$ and $f_{xyy}$.
\end{exercise}
\begin{solbox}
$f_x=2xe^y+y\cos x$.\quad $f_{xx}=2e^y-y\sin x$.
\[
  f_{xxy}=\pd{}{y}(2e^y-y\sin x)=2e^y-\sin x.
\]
$f_{xy}=\pd{}{y}(2xe^y+y\cos x)=2xe^y+\cos x$.
\[
  f_{xyy}=\pd{}{y}(2xe^y+\cos x)=2xe^y.
\]
\end{solbox}

% ============================================================
\section{Tangent Planes and Linear Approximations}
% ============================================================

%------------------------------------------------------------------
\begin{exercise}
(a) Tangent plane to $z=x^2+2y^2-3x$ at $(1,1,0)$;
(b) tangent plane to $z=e^{x-y}$ at $(2,2,1)$.
\end{exercise}
\begin{solbox}
\textbf{(a)} $f_x=2x-3$, $f_y=4y$. At $(1,1)$: $f_x=-1$, $f_y=4$.
\[
  z-0=-1(x-1)+4(y-1)\implies \boxed{z=-x+4y-3}.
\]

\textbf{(b)} $f_x=e^{x-y}$, $f_y=-e^{x-y}$. At $(2,2)$: $f_x=1$, $f_y=-1$.
\[
  z-1=1(x-2)+(-1)(y-2)\implies \boxed{z=x-y+1}.
\]
\end{solbox}

%------------------------------------------------------------------
\begin{exercise}
Find the linearisation of $f(x,y)=\sqrt{20-x^2-7y^2}$ at $(2,1)$ and estimate $f(1.95,1.05)$.
\end{exercise}
\begin{solbox}
$f(2,1)=\sqrt{20-4-7}=3$.
\[
  f_x=\frac{-x}{\sqrt{20-x^2-7y^2}}\Bigg|_{(2,1)}=\frac{-2}{3},\qquad
  f_y=\frac{-7y}{\sqrt{20-x^2-7y^2}}\Bigg|_{(2,1)}=\frac{-7}{3}.
\]
\[
  L(x,y)=3-\tfrac{2}{3}(x-2)-\tfrac{7}{3}(y-1).
\]
\[
  f(1.95,1.05)\approx L(1.95,1.05)=3-\tfrac{2}{3}(-0.05)-\tfrac{7}{3}(0.05)
  =3+\tfrac{0.10}{3}-\tfrac{0.35}{3}=3-\tfrac{0.25}{3}\approx 2.917.
\]
\end{solbox}

%------------------------------------------------------------------
\begin{exercise}
Box $80\times60\times50\,\text{cm}$, error $\pm0.2\,\text{cm}$ in each dimension.
Estimate maximum error in volume.
\end{exercise}
\begin{solbox}
$V=lwh$, $\dd V=wh\,\dd l+lh\,\dd w+lw\,\dd h$.
Maximum error:
\[
  |dV|\le(wh+lh+lw)\cdot 0.2=(60\cdot50+80\cdot50+80\cdot60)\cdot0.2
  =(3000+4000+4800)\cdot0.2=11800\cdot0.2=\boxed{2360\,\text{cm}^3}.
\]
\end{solbox}

%------------------------------------------------------------------
\begin{exercise}
Find the total differential $dz$ for (a) $z=x^4y+x\sqrt{y}$;
(b) $z=\arctan\!\bigl(\tfrac{x}{y}\bigr)$.
\end{exercise}
\begin{solbox}
\textbf{(a)}
$\pd{z}{x}=4x^3y+\sqrt{y}$,\quad $\pd{z}{y}=x^4+\dfrac{x}{2\sqrt{y}}$.
\[
  dz=\bigl(4x^3y+\sqrt{y}\bigr)\dif{x}+\biggl(x^4+\frac{x}{2\sqrt{y}}\biggr)\dif{y}.
\]

\textbf{(b)} $\pd{z}{x}=\dfrac{1}{1+x^2/y^2}\cdot\dfrac{1}{y}=\dfrac{y}{x^2+y^2}$,\quad
$\pd{z}{y}=\dfrac{1}{1+x^2/y^2}\cdot\dfrac{-x}{y^2}=\dfrac{-x}{x^2+y^2}$.
\[
  dz=\frac{y}{x^2+y^2}\dif{x}-\frac{x}{x^2+y^2}\dif{y}.
\]
\end{solbox}

% ============================================================
\section{The Chain Rule}
% ============================================================

%------------------------------------------------------------------
\begin{exercise}
Use the Chain Rule to find $dz/dt$.
\end{exercise}
\begin{solbox}
\textbf{(a)} $z=x^2y+3xy^4$, $x=\sin t$, $y=e^t$.
\[
  \frac{dz}{dt}=\underbrace{(2xy+3y^4)}_{\partial z/\partial x}\cos t
               +\underbrace{(x^2+12xy^3)}_{\partial z/\partial y}e^t.
\]
(General formula; substituting $x=\sin t$, $y=e^t$ gives the explicit form.)

\textbf{(b)} $z=\sqrt{x^2+y^2}$, $x=e^t\cos t$, $y=e^t\sin t$.
Note $x^2+y^2=e^{2t}$, so $z=e^t$.
\[
  \frac{dz}{dt}=\frac{x}{e^t}(e^t\cos t-e^t\sin t)+\frac{y}{e^t}(e^t\sin t+e^t\cos t).
\]
\[
  =e^t\cos t(\cos t-\sin t)+e^t\sin t(\sin t+\cos t)=e^t(\cos^2\!t+\sin^2\!t)=\boxed{e^t}.
\]
\end{solbox}

%------------------------------------------------------------------
\begin{exercise}
Use the Chain Rule to find $\partial z/\partial s$ and $\partial z/\partial t$.
\end{exercise}
\begin{solbox}
\textbf{(a)} $z=e^{2x+3y}$, $x=s/t$, $y=t/s$.
\[
  \pd{z}{s}=e^{2x+3y}\!\left(2\cdot\frac{1}{t}+3\cdot\frac{-t}{s^2}\right)
  =e^{2s/t+3t/s}\!\left(\frac{2}{t}-\frac{3t}{s^2}\right).
\]
\[
  \pd{z}{t}=e^{2x+3y}\!\left(2\cdot\frac{-s}{t^2}+3\cdot\frac{1}{s}\right)
  =e^{2s/t+3t/s}\!\left(\frac{3}{s}-\frac{2s}{t^2}\right).
\]

\textbf{(b)} $z=x^2\tan y$, $x=st^2$, $y=s^2t$.
\begin{align*}
  \pd{z}{s}&=2x\tan y\cdot t^2+x^2\sec^2\!y\cdot 2st
  =2st^4\tan(s^2t)+2s^3t^5\sec^2(s^2t),\\
  \pd{z}{t}&=2x\tan y\cdot 2st+x^2\sec^2\!y\cdot s^2
  =4s^2t^3\tan(s^2t)+s^4t^4\sec^2(s^2t).
\end{align*}
\end{solbox}

%------------------------------------------------------------------
\begin{exercise}[Implicit Differentiation]
(a) Find $dy/dx$ if $x^4+x^2y^2+y^4=3$;
(b) find $\partial z/\partial x$ and $\partial z/\partial y$ if $x^2+2y^2+3z^2+xyz=6$.
\end{exercise}
\begin{solbox}
\textbf{(a)} Let $F=x^4+x^2y^2+y^4-3$.
$F_x=4x^3+2xy^2$, $F_y=2x^2y+4y^3$.
\[
  \frac{dy}{dx}=-\frac{F_x}{F_y}=-\frac{4x^3+2xy^2}{2x^2y+4y^3}=-\frac{2x^3+xy^2}{x^2y+2y^3}.
\]

\textbf{(b)} Let $F=x^2+2y^2+3z^2+xyz-6$.
$F_x=2x+yz$, $F_y=4y+xz$, $F_z=6z+xy$.
\[
  \pd{z}{x}=-\frac{F_x}{F_z}=-\frac{2x+yz}{6z+xy},\qquad
  \pd{z}{y}=-\frac{F_y}{F_z}=-\frac{4y+xz}{6z+xy}.
\]
\end{solbox}

% ============================================================
\section{Directional Derivatives and the Gradient}
% ============================================================

%------------------------------------------------------------------
\begin{exercise}
Find the directional derivative at the given point in the given direction.
\end{exercise}
\begin{solbox}
\textbf{(a)} $f=x^2y^3-2x$, point $(1,-1)$, direction angle $\theta=\pi/4$.
Unit vector $\mathbf{u}=\langle\tfrac{1}{\sqrt{2}},\tfrac{1}{\sqrt{2}}\rangle$.
$f_x=2xy^3-2$, $f_y=3x^2y^2$. At $(1,-1)$: $f_x=-4$, $f_y=3$.
\[
  D_{\mathbf{u}}f=(-4)\frac{1}{\sqrt{2}}+3\frac{1}{\sqrt{2}}=\frac{-1}{\sqrt{2}}=-\frac{\sqrt{2}}{2}.
\]

\textbf{(b)} $f=x^2yz+xz^2$, point $(1,1,1)$, $\mathbf{v}=\langle1,2,-1\rangle$.
$|\mathbf{v}|=\sqrt{6}$, $\mathbf{u}=\langle1,2,-1\rangle/\sqrt{6}$.
$f_x=2xyz+z^2=3$, $f_y=x^2z=1$, $f_z=x^2y+2xz=3$ at $(1,1,1)$.
\[
  D_{\mathbf{u}}f=\frac{3(1)+1(2)+3(-1)}{\sqrt{6}}=\frac{2}{\sqrt{6}}=\frac{\sqrt{6}}{3}.
\]
\end{solbox}

%------------------------------------------------------------------
\begin{exercise}
(a) Gradient of $f=\sin(x+2y)$ at $(\pi/6,0)$;
(b) maximum rate of change of $f=\sqrt{x^2+y^2+z^2}$ at $(1,2,2)$.
\end{exercise}
\begin{solbox}
\textbf{(a)} $\grad f=\langle\cos(x+2y),\;2\cos(x+2y)\rangle$.
At $(\pi/6,0)$: $x+2y=\pi/6$, so $\grad f=\bigl\langle\tfrac{\sqrt{3}}{2},\sqrt{3}\bigr\rangle$.

\textbf{(b)} $\grad f=\dfrac{\langle x,y,z\rangle}{\sqrt{x^2+y^2+z^2}}$.
At $(1,2,2)$: $|\grad f|=\dfrac{\sqrt{1+4+4}}{3}=1$.
Maximum rate of change $=1$, in direction $\langle1,2,2\rangle/3$.
\end{solbox}

%------------------------------------------------------------------
\begin{exercise}
$T(x,y)=200e^{-(x^2+y^2)/100}$ at $(3,4)$:
(a) direction of fastest warming; (b) rate in that direction;
(c) rate in direction $\mathbf{u}=\langle4,3\rangle/5$.
\end{exercise}
\begin{solbox}
$\grad T=200e^{-(x^2+y^2)/100}\cdot\left\langle\frac{-2x}{100},\frac{-2y}{100}\right\rangle
=-4e^{-(x^2+y^2)/100}\langle x,y\rangle$.

At $(3,4)$: $x^2+y^2=25$, $e^{-1/4}$.
$\grad T(3,4)=-4e^{-1/4}\langle3,4\rangle$.

\textbf{(a)} Fastest warming: in direction $-\langle3,4\rangle/5$ (toward origin, since $\grad T$ points inward).

\textbf{(b)} $|\grad T|=4e^{-1/4}\cdot5=20e^{-1/4}\approx15.6$ per unit length.

\textbf{(c)} $D_{\mathbf{u}}T=\grad T\cdot\mathbf{u}=(-4e^{-1/4})\langle3,4\rangle\cdot\frac{\langle4,3\rangle}{5}
=\frac{-4e^{-1/4}(12+12)}{5}=-\frac{96e^{-1/4}}{5}\approx-14.9$.
\end{solbox}

% ============================================================
\section{Extrema of Functions of Several Variables}
% ============================================================

%------------------------------------------------------------------
\begin{exercise}
Find and classify all critical points of each function.
\end{exercise}
\begin{solbox}
\textbf{(a)} $f=x^2+xy+y^2-3x$. Setting $f_x=2x+y-3=0$ and $f_y=x+2y=0$:
from $f_y$, $x=-2y$; substituting: $-4y+y-3=0\Rightarrow y=1$, $x=-2$.
$D=f_{xx}f_{yy}-f_{xy}^2=2\cdot2-1=3>0$, $f_{xx}=2>0$: \textbf{local min} at $(-2,1)$, value $f(-2,1)=4-2+1+6=9$.

\textbf{(b)} $f=4xy-x^4-y^4$. $f_x=4y-4x^3=0$, $f_y=4x-4y^3=0$.
From these: $y=x^3$ and $x=y^3$, giving $x=x^9\Rightarrow x=0,\pm1$.
\begin{itemize}
  \item $(0,0)$: $D=(-12x^2)(-12y^2)-16\big|_{(0,0)}=-16<0$: \textbf{saddle point}.
  \item $(1,1)$: $D=144-16=128>0$, $f_{xx}=-12<0$: \textbf{local max}, $f(1,1)=2$.
  \item $(-1,-1)$: $D=128>0$, $f_{xx}=-12<0$: \textbf{local max}, $f(-1,-1)=2$.
\end{itemize}

\textbf{(c)} $f=e^x(x+y^2+2y)$. Since $e^x>0$:
$f_x=e^x(x+1+y^2+2y)=0\Rightarrow x=-(y+1)^2$.\quad $f_y=e^x(2y+2)=0\Rightarrow y=-1$.
At $y=-1$: $x=0$. Critical point: $(0,-1)$.
$f_{xx}=e^x(x+2+y^2+2y)$, $f_{yy}=2e^x$, $f_{xy}=e^x(2y+2)$.
At $(0,-1)$: $f_{xx}=1$, $f_{yy}=2$, $f_{xy}=0$.
$D=2>0$, $f_{xx}>0$: \textbf{local min} at $(0,-1)$, value $f(0,-1)=-1$.
\end{solbox}

%------------------------------------------------------------------
\begin{exercise}
Absolute max/min of $f(x,y)=x^2-2xy+2y$ on the triangle with vertices $(0,0)$, $(3,0)$, $(3,3)$.
\end{exercise}
\begin{solbox}
\emph{Interior:} $f_x=2x-2y=0$, $f_y=-2x+2=0\Rightarrow x=1$, $y=1$. $f(1,1)=1$.

\emph{Boundary:}
\begin{itemize}
  \item $y=0$, $0\le x\le3$: $f=x^2$. Min $0$ at $x=0$; max $9$ at $x=3$.
  \item $x=3$, $0\le y\le3$: $f=9-6y+2y=9-4y$. Max $9$ at $y=0$; min $-3$ at $y=3$.
  \item $y=x$, $0\le x\le3$: $f=-x^2+2x$. $f'=-2x+2=0\Rightarrow x=1$, $f(1)=1$. Endpoints: $0$ and $-3$.
\end{itemize}
\textbf{Absolute max} $=9$ at $(3,0)$.\quad \textbf{Absolute min} $=-3$ at $(3,3)$.
\end{solbox}

%------------------------------------------------------------------
\begin{exercise}
Find three positive numbers $a,b,c$ with $a+b+c=100$ maximising $abc$.
\end{exercise}
\begin{solbox}
By the AM-GM inequality: $\dfrac{a+b+c}{3}\ge(abc)^{1/3}$, with equality iff $a=b=c$.
Setting $a=b=c=\dfrac{100}{3}$ gives the maximum product
$\left(\dfrac{100}{3}\right)^3=\dfrac{10^6}{27}\approx37037$.
\end{solbox}

% ============================================================
\section{Lagrange Multipliers}
% ============================================================

%------------------------------------------------------------------
\begin{exercise}
Use Lagrange multipliers.
\end{exercise}
\begin{solbox}
\textbf{(a)} Maximise $f=xy$ subject to $g=x+2y-8=0$.
$\grad f=\lambda\grad g$: $(y,x)=\lambda(1,2)$, so $y=\lambda$ and $x=2\lambda\Rightarrow x=2y$.
Substituting: $2y+2y=8\Rightarrow y=2,\,x=4$. Maximum $f(4,2)=8$.

\textbf{(b)} Minimise $f=x^2+y^2+z^2$ subject to $x+y+z=12$.
$(2x,2y,2z)=\lambda(1,1,1)\Rightarrow x=y=z=\lambda/2$. Sum: $3\lambda/2=12\Rightarrow\lambda=8$, $x=y=z=4$.
Minimum $f(4,4,4)=48$.

\textbf{(c)} Extrema of $f=x^2+2y^2$ on $x^2+y^2=1$.
$(2x,4y)=\lambda(2x,2y)$.
\begin{itemize}
  \item $x\neq0$: $\lambda=1$, then $4y=2y\Rightarrow y=0$, $x=\pm1$: $f=1$.
  \item $x=0$: $y=\pm1$: $f=2$.
\end{itemize}
\textbf{Minimum} $=1$ at $(\pm1,0)$.\quad\textbf{Maximum} $=2$ at $(0,\pm1)$.
\end{solbox}

%------------------------------------------------------------------
\begin{exercise}
Nearest point on the plane $2x+y+z=5$ to the origin.
\end{exercise}
\begin{solbox}
Minimise $f=x^2+y^2+z^2$ subject to $g=2x+y+z-5=0$.
$(2x,2y,2z)=\lambda(2,1,1)\Rightarrow x=\lambda,\,y=\tfrac{\lambda}{2},\,z=\tfrac{\lambda}{2}$.
Substituting: $2\lambda+\tfrac{\lambda}{2}+\tfrac{\lambda}{2}=5\Rightarrow3\lambda=5\Rightarrow\lambda=\tfrac{5}{3}$.
\[
  \text{Nearest point: }\left(\tfrac{5}{3},\,\tfrac{5}{6},\,\tfrac{5}{6}\right).
\]
\end{solbox}

%------------------------------------------------------------------
\begin{exercise}
Open-top rectangular box from $12\,\text{m}^2$ of cardboard; maximise volume.
\end{exercise}
\begin{solbox}
Dimensions: $x,y,z$ (height $z$). Surface: $g=xy+2xz+2yz=12$.
Maximise $V=xyz$.
$\grad V=\lambda\grad g$:
\[
  yz=\lambda(y+2z),\quad xz=\lambda(x+2z),\quad xy=\lambda(2x+2y).
\]
From the first two: $yz(x+2z)=xz(y+2z)\Rightarrow xyz+2yz^2=xyz+2xz^2\Rightarrow y=x$.
Substituting $y=x$: $xz=\lambda(x+2z)$ and $x^2=4\lambda x\Rightarrow\lambda=x/4$.
Then $xz=(x/4)(x+2z)\Rightarrow4xz=x^2+2xz\Rightarrow2xz=x^2\Rightarrow z=x/2$.
Surface: $x^2+2x(x/2)+2x(x/2)=12\Rightarrow3x^2=12\Rightarrow x=2$.
\[
  \text{Optimal dimensions: }x=y=2\,\text{m},\;z=1\,\text{m};\quad V_{\max}=4\,\text{m}^3.
\]
\end{solbox}

% ============================================================
\section{Double Integrals over Rectangles}
% ============================================================

%------------------------------------------------------------------
\begin{exercise}
Evaluate the iterated integrals.
\end{exercise}
\begin{solbox}
\textbf{(a)} $\displaystyle\int_0^3\int_1^2(x^2y-2xy^3)\dif{y}\dif{x}$.
\[
  \int_1^2(x^2y-2xy^3)\dif{y}=\left[\frac{x^2y^2}{2}-\frac{xy^4}{2}\right]_1^2
  =2x^2-8x-\frac{x^2}{2}+\frac{x}{2}=\frac{3x^2}{2}-\frac{15x}{2}.
\]
\[
  \int_0^3\left(\frac{3x^2}{2}-\frac{15x}{2}\right)\dif{x}
  =\left[\frac{x^3}{2}-\frac{15x^2}{4}\right]_0^3=\frac{27}{2}-\frac{135}{4}=-\frac{81}{4}.
\]

\textbf{(b)} $\displaystyle\int_0^\pi\int_0^1 y\sin(xy)\dif{y}\dif{x}$.
\[
  \int_0^1 y\sin(xy)\dif{y}=\left[-\cos(xy)\right]_0^1=1-\cos x.
\]
\[
  \int_0^\pi(1-\cos x)\dif{x}=\bigl[x-\sin x\bigr]_0^\pi=\pi.
\]

\textbf{(c)} $\displaystyle\int_0^1\int_0^1 xe^{xy}\dif{y}\dif{x}$.
\[
  \int_0^1 xe^{xy}\dif{y}=\bigl[e^{xy}\bigr]_0^1=e^x-1.
\]
\[
  \int_0^1(e^x-1)\dif{x}=\bigl[e^x-x\bigr]_0^1=e-1-1=e-2.
\]
\end{solbox}

%------------------------------------------------------------------
\begin{exercise}
Evaluate $\iint_R(x+2y)\dif{A}$, $R=[0,2]\times[0,3]$ in both orders.
\end{exercise}
\begin{solbox}
\textbf{Order $dy\,dx$:}
$\int_0^2\int_0^3(x+2y)\dif{y}\dif{x}=\int_0^2\bigl[xy+y^2\bigr]_0^3\dif{x}=\int_0^2(3x+9)\dif{x}=\bigl[\tfrac{3x^2}{2}+9x\bigr]_0^2=6+18=24$.

\textbf{Order $dx\,dy$:}
$\int_0^3\int_0^2(x+2y)\dif{x}\dif{y}=\int_0^3\bigl[\tfrac{x^2}{2}+2xy\bigr]_0^2\dif{y}=\int_0^3(2+4y)\dif{y}=\bigl[2y+2y^2\bigr]_0^3=6+18=24$.

Both orders give $\boxed{24}$.\hfill$\checkmark$
\end{solbox}

%------------------------------------------------------------------
\begin{exercise}
Volume above $R=[-1,1]\times[0,2]$ and below $z=2+x+3y$.
\end{exercise}
\begin{solbox}
\[
  V=\int_{-1}^1\int_0^2(2+x+3y)\dif{y}\dif{x}
  =\int_{-1}^1\bigl[2y+xy+\tfrac{3y^2}{2}\bigr]_0^2\dif{x}
  =\int_{-1}^1(4+2x+6)\dif{x}=\int_{-1}^1(10+2x)\dif{x}
\]
\[
  =\bigl[10x+x^2\bigr]_{-1}^1=(10+1)-(-10+1)=\boxed{20}.
\]
\end{solbox}

% ============================================================
\section{Double Integrals over General Regions}
% ============================================================

%------------------------------------------------------------------
\begin{exercise}
Evaluate $\iint_D f\,\dif{A}$ over each described region.
\end{exercise}
\begin{solbox}
\textbf{(a)} $D$: bounded by $y=x^2$ and $y=x+2$. Intersections: $x=-1,2$.
\begin{align*}
  &\int_{-1}^2\int_{x^2}^{x+2}(2x-y)\dif{y}\dif{x}
  =\int_{-1}^2\left[2xy-\frac{y^2}{2}\right]_{x^2}^{x+2}\dif{x}\\
  &=\int_{-1}^2\!\Bigl(2x(x+2)-\tfrac{(x+2)^2}{2}-2x^3+\tfrac{x^4}{2}\Bigr)\dif{x}
  =\int_{-1}^2\!\Bigl(\tfrac{x^4}{2}-2x^3+\tfrac{3x^2}{2}+2x-2\Bigr)\dif{x}\\
  &=\left[\frac{x^5}{10}-\frac{x^4}{2}+\frac{x^3}{2}+x^2-2x\right]_{-1}^2
  =\left(\frac{32}{10}-8+4+4-4\right)-\left(\frac{-1}{10}-\frac{1}{2}-\frac{1}{2}+1+2\right)\\
  &=\frac{32}{10}-4-\left(\frac{-1}{10}+2\right)
  =\frac{33}{10}-6=-\frac{27}{10}.
\end{align*}

\textbf{(b)} $D$: bounded by $y=x$ and $y=x^2$, $0\le x\le1$.
\[
  \int_0^1\int_{x^2}^x xy\dif{y}\dif{x}=\int_0^1 x\cdot\frac{x^2-x^4}{2}\dif{x}
  =\frac{1}{2}\int_0^1(x^3-x^5)\dif{x}=\frac{1}{2}\left[\frac{1}{4}-\frac{1}{6}\right]=\frac{1}{24}.
\]

\textbf{(c)} $D$: triangle $(0,0),(1,0),(0,1)$; $0\le x\le1$, $0\le y\le1-x$.
\[
  \int_0^1\int_0^{1-x}x\dif{y}\dif{x}=\int_0^1 x(1-x)\dif{x}=\left[\frac{x^2}{2}-\frac{x^3}{3}\right]_0^1=\frac{1}{6}.
\]
\end{solbox}

%------------------------------------------------------------------
\begin{exercise}
Reverse the order and evaluate.
\end{exercise}
\begin{solbox}
\textbf{(a)} $\int_0^1\int_x^1\sin(y^2)\dif{y}\dif{x}$.
Region: $0\le y\le1$, $0\le x\le y$.
\[
  =\int_0^1\int_0^y\sin(y^2)\dif{x}\dif{y}=\int_0^1 y\sin(y^2)\dif{y}
  =\left[-\frac{\cos(y^2)}{2}\right]_0^1=\frac{1-\cos1}{2}.
\]

\textbf{(b)} $\int_0^4\int_{\sqrt{y}}^2 e^{x^3}\dif{x}\dif{y}$.
Region: $0\le x\le2$, $0\le y\le x^2$.
\[
  =\int_0^2\int_0^{x^2}e^{x^3}\dif{y}\dif{x}=\int_0^2 x^2e^{x^3}\dif{x}
  =\left[\frac{e^{x^3}}{3}\right]_0^2=\frac{e^8-1}{3}.
\]
\end{solbox}

%------------------------------------------------------------------
\begin{exercise}
Area between $y=3x-x^2$ and $y=x$.
\end{exercise}
\begin{solbox}
Intersections: $3x-x^2=x\Rightarrow 2x-x^2=0\Rightarrow x=0,2$.
\[
  \text{Area}=\int_0^2\bigl[(3x-x^2)-x\bigr]\dif{x}=\int_0^2(2x-x^2)\dif{x}
  =\left[x^2-\frac{x^3}{3}\right]_0^2=4-\frac{8}{3}=\frac{4}{3}.
\]
\end{solbox}

% ============================================================
\section{Double Integrals in Polar Coordinates}
% ============================================================

%------------------------------------------------------------------
\begin{exercise}
Convert to polar and evaluate.
\end{exercise}
\begin{solbox}
\textbf{(a)} $\iint_D\sqrt{x^2+y^2}\,\dif{A}$, $D$: right half-disk of radius 2.
$r\in[0,2]$, $\theta\in[-\pi/2,\pi/2]$.
\[
  =\int_{-\pi/2}^{\pi/2}\!\int_0^2 r\cdot r\dif{r}\dif{\theta}
  =\pi\cdot\left[\frac{r^3}{3}\right]_0^2=\frac{8\pi}{3}.
\]

\textbf{(b)} $\iint_D e^{-(x^2+y^2)}\dif{A}$, $D$: disk of radius $a$.
\[
  =\int_0^{2\pi}\!\int_0^a e^{-r^2}r\dif{r}\dif{\theta}
  =2\pi\left[-\frac{e^{-r^2}}{2}\right]_0^a=\pi\bigl(1-e^{-a^2}\bigr).
\]

\textbf{(c)} $\int_0^2\!\int_0^{\sqrt{4-x^2}}(x^2+y^2)^{3/2}\dif{y}\dif{x}$.
Upper semicircle: $r\in[0,2]$, $\theta\in[0,\pi]$.
\[
  =\int_0^\pi\!\int_0^2 r^3\cdot r\dif{r}\dif{\theta}
  =\pi\left[\frac{r^5}{5}\right]_0^2=\frac{32\pi}{5}.
\]
\end{solbox}

%------------------------------------------------------------------
\begin{exercise}
Volume under $z=9-x^2-y^2$, above $xy$-plane, inside $x^2+y^2=4$.
\end{exercise}
\begin{solbox}
\[
  V=\int_0^{2\pi}\!\int_0^2(9-r^2)r\dif{r}\dif{\theta}
  =2\pi\int_0^2(9r-r^3)\dif{r}
  =2\pi\left[\frac{9r^2}{2}-\frac{r^4}{4}\right]_0^2=2\pi(18-4)=\boxed{28\pi}.
\]
\end{solbox}

%------------------------------------------------------------------
\begin{exercise}
Area inside $r=1+\cos\theta$ and outside $r=1$.
\end{exercise}
\begin{solbox}
Intersection: $\cos\theta=0\Rightarrow\theta=\pm\pi/2$.
\[
  A=\frac{1}{2}\int_{-\pi/2}^{\pi/2}\!\bigl[(1+\cos\theta)^2-1\bigr]\dif{\theta}
  =\frac{1}{2}\int_{-\pi/2}^{\pi/2}(2\cos\theta+\cos^2\!\theta)\dif{\theta}.
\]
\[
  =\frac{1}{2}\left[2\sin\theta+\frac{\theta}{2}+\frac{\sin2\theta}{4}\right]_{-\pi/2}^{\pi/2}
  =\frac{1}{2}\!\left[\left(2+\frac{\pi}{4}\right)-\left(-2-\frac{\pi}{4}\right)\right]
  =\frac{1}{2}\!\left(4+\frac{\pi}{2}\right)=2+\frac{\pi}{4}.
\]
\end{solbox}

% ============================================================
\section{Triple Integrals}
% ============================================================

%------------------------------------------------------------------
\begin{exercise}
Evaluate the triple integrals.
\end{exercise}
\begin{solbox}
\textbf{(a)} $\int_0^2\!\int_0^x\!\int_0^{x+y}6xy\dif{z}\dif{y}\dif{x}$.
\[
  \int_0^{x+y}6xy\dif{z}=6xy(x+y).
\]
\[
  \int_0^x6xy(x+y)\dif{y}=6x\!\int_0^x(xy+y^2)\dif{y}=6x\!\left[\frac{xy^2}{2}+\frac{y^3}{3}\right]_0^x=6x\!\left(\frac{x^3}{2}+\frac{x^3}{3}\right)=5x^5.
\]
\[
  \int_0^25x^5\dif{x}=\left[x^5\right]_0^2=32.
\]

\textbf{(b)} $\iiint_B xe^{y+2z}\dif{V}$, $B=[0,1]^3$. Separable:
\[
  =\left(\int_0^1 x\dif{x}\right)\!\left(\int_0^1 e^y\dif{y}\right)\!\left(\int_0^1 e^{2z}\dif{z}\right)
  =\frac{1}{2}\cdot(e-1)\cdot\frac{e^2-1}{2}=\frac{(e-1)(e^2-1)}{4}=\frac{(e-1)^2(e+1)}{4}.
\]

\textbf{(c)} $\iiint_E y\dif{V}$, $0\le x\le1$, $0\le y\le1$, $0\le z\le x+y+5$.
\[
  =\int_0^1\!\int_0^1 y(x+y+5)\dif{y}\dif{x}
  =\int_0^1\!\left[\frac{xy^2}{2}+\frac{y^3}{3}+\frac{5y^2}{2}\right]_0^1\dif{x}
  =\int_0^1\!\left(\frac{x}{2}+\frac{1}{3}+\frac{5}{2}\right)\dif{x}
  =\left[\frac{x^2}{4}+\frac{17x}{6}\right]_0^1=\frac{1}{4}+\frac{17}{6}=\frac{37}{12}.
\]
\end{solbox}

%------------------------------------------------------------------
\begin{exercise}
(a) $\iiint_E z\dif{V}$ for tetrahedron; (b) volume between two paraboloids.
\end{exercise}
\begin{solbox}
\textbf{(a)} Tetrahedron: $0\le x\le1$, $0\le y\le1-x$, $0\le z\le1-x-y$.
\begin{align*}
  &\int_0^1\!\int_0^{1-x}\!\int_0^{1-x-y}z\dif{z}\dif{y}\dif{x}
  =\int_0^1\!\int_0^{1-x}\frac{(1-x-y)^2}{2}\dif{y}\dif{x}\\
  &=\frac{1}{2}\int_0^1\!\left[-\frac{(1-x-y)^3}{3}\right]_0^{1-x}\dif{x}
  =\frac{1}{6}\int_0^1(1-x)^3\dif{x}
  =\frac{1}{6}\left[-\frac{(1-x)^4}{4}\right]_0^1=\frac{1}{24}.
\end{align*}

\textbf{(b)} Intersection of $z=x^2+y^2$ and $z=36-3(x^2+y^2)$:
$4(x^2+y^2)=36\Rightarrow r=3$.
\[
  V=\int_0^{2\pi}\!\int_0^3\bigl[(36-3r^2)-r^2\bigr]r\dif{r}\dif{\theta}
  =2\pi\int_0^3(36r-4r^3)\dif{r}=2\pi\bigl[18r^2-r^4\bigr]_0^3=2\pi(162-81)=\boxed{162\pi}.
\]
\end{solbox}

% ============================================================
\section{Triple Integrals in Cylindrical Coordinates}
% ============================================================

%------------------------------------------------------------------
\begin{exercise}
Use cylindrical coordinates to evaluate.
\end{exercise}
\begin{solbox}
\textbf{(a)} $\iiint_E(x^2+y^2)\dif{V}$, inside $r=2$, $-3\le z\le5$.
\[
  =\int_0^{2\pi}\!\int_0^2\!\int_{-3}^5 r^2\cdot r\dif{z}\dif{r}\dif{\theta}
  =2\pi\cdot8\int_0^2 r^3\dif{r}=16\pi\left[\frac{r^4}{4}\right]_0^2=16\pi\cdot4=\boxed{64\pi}.
\]

\textbf{(b)} $\iiint_E\sqrt{x^2+y^2}\dif{V}$, below $z=4$, above $z=r^2$.
\[
  =\int_0^{2\pi}\!\int_0^2\!\int_{r^2}^4 r\cdot r\dif{z}\dif{r}\dif{\theta}
  =2\pi\int_0^2 r^2(4-r^2)\dif{r}=2\pi\int_0^2(4r^2-r^4)\dif{r}
  =2\pi\left[\frac{4r^3}{3}-\frac{r^5}{5}\right]_0^2=2\pi\!\left(\frac{32}{3}-\frac{32}{5}\right)=\frac{128\pi}{15}.
\]

\textbf{(c)} Convert: $r\in[0,2]$, $\theta\in[0,2\pi]$, $z\in[r,2]$.
\[
  =\int_0^{2\pi}\!\int_0^2\!\int_r^2 r^2\cdot r\dif{z}\dif{r}\dif{\theta}
  =2\pi\int_0^2 r^3(2-r)\dif{r}=2\pi\!\int_0^2(2r^3-r^4)\dif{r}
  =2\pi\!\left[\frac{r^4}{2}-\frac{r^5}{5}\right]_0^2=2\pi\!\left(8-\frac{32}{5}\right)=\frac{16\pi}{5}.
\]
\end{solbox}

%------------------------------------------------------------------
\begin{exercise}
Volume of upper hemisphere of radius $a$ using cylindrical coordinates.
\end{exercise}
\begin{solbox}
\[
  V=\int_0^{2\pi}\!\int_0^a\!\int_0^{\sqrt{a^2-r^2}}r\dif{z}\dif{r}\dif{\theta}
  =2\pi\int_0^a r\sqrt{a^2-r^2}\dif{r}.
\]
Let $u=a^2-r^2$, $\dif{u}=-2r\dif{r}$:
\[
  =2\pi\cdot\frac{1}{2}\int_0^{a^2}\sqrt{u}\dif{u}=\pi\!\left[\frac{2u^{3/2}}{3}\right]_0^{a^2}=\frac{2\pi a^3}{3}.
\]
\end{solbox}

% ============================================================
\section{Triple Integrals in Spherical Coordinates}
% ============================================================

%------------------------------------------------------------------
\begin{exercise}
Use spherical coordinates to evaluate.
\end{exercise}
\begin{solbox}
\textbf{(a)} $\iiint_B e^{(\rho^2)^{3/2}}\rho^2\sin\varphi\dif{\rho}\dif{\varphi}\dif{\theta}$
with $\rho\in[0,1]$.
\[
  =\int_0^{2\pi}\dif\theta\int_0^\pi\sin\varphi\dif\varphi\int_0^1 e^{\rho^3}\rho^2\dif\rho
  =2\pi\cdot2\cdot\left[\frac{e^{\rho^3}}{3}\right]_0^1=4\pi\cdot\frac{e-1}{3}=\frac{4\pi(e-1)}{3}.
\]

\textbf{(b)} $\iiint_H(x^2+y^2+z^2)^2\dif{V}$, upper hemisphere $\rho\le2$, $0\le\varphi\le\pi/2$.
\[
  =\int_0^{2\pi}\dif\theta\int_0^{\pi/2}\sin\varphi\dif\varphi\int_0^2\rho^4\cdot\rho^2\dif\rho
  =2\pi\cdot[-\cos\varphi]_0^{\pi/2}\cdot\left[\frac{\rho^7}{7}\right]_0^2
  =2\pi\cdot1\cdot\frac{128}{7}=\frac{256\pi}{7}.
\]

\textbf{(c)} $\iiint_E z\dif{V}$, above cone $\varphi=\pi/4$, inside sphere $\rho=2$.
\[
  =\int_0^{2\pi}\dif\theta\int_0^{\pi/4}\int_0^2\rho\cos\varphi\cdot\rho^2\sin\varphi\dif\rho\dif\varphi
  =2\pi\left[\frac{\rho^4}{4}\right]_0^2\int_0^{\pi/4}\sin\varphi\cos\varphi\dif\varphi
  =2\pi\cdot4\cdot\left[\frac{\sin^2\!\varphi}{2}\right]_0^{\pi/4}
  =8\pi\cdot\frac{1}{4}=2\pi.
\]
\end{solbox}

%------------------------------------------------------------------
\begin{exercise}
(a) Volume inside sphere $\rho=2$, above cone $\varphi=\pi/4$;
(b) $\iiint_E(x^2+y^2)\dif{V}$ for ball of radius 3.
\end{exercise}
\begin{solbox}
\textbf{(a)} $\varphi\in[0,\pi/4]$, $\rho\in[0,2]$:
\[
  V=\int_0^{2\pi}\!\int_0^{\pi/4}\!\int_0^2\rho^2\sin\varphi\dif\rho\dif\varphi\dif\theta
  =2\pi\cdot\frac{8}{3}\cdot[-\cos\varphi]_0^{\pi/4}
  =\frac{16\pi}{3}\!\left(1-\frac{1}{\sqrt{2}}\right)=\frac{8\pi(2-\sqrt{2})}{3}.
\]

\textbf{(b)} $x^2+y^2=\rho^2\sin^2\!\varphi$:
\[
  \iiint_E(x^2+y^2)\dif{V}=\int_0^{2\pi}\!\int_0^\pi\!\int_0^3\rho^2\sin^2\!\varphi\cdot\rho^2\sin\varphi\dif\rho\dif\varphi\dif\theta
  =2\pi\cdot\frac{3^5}{5}\cdot\int_0^\pi\sin^3\!\varphi\dif\varphi.
\]
$\int_0^\pi\sin^3\!\varphi\dif\varphi=\tfrac{4}{3}$.
\[
  =2\pi\cdot\frac{243}{5}\cdot\frac{4}{3}=\frac{648\pi}{5}.
\]
\end{solbox}

% ============================================================
\section{Mixed Challenge Problems}
% ============================================================

%------------------------------------------------------------------
\begin{exercise}
$T(x,y)=x^2-xy+y^2$. (a) Critical points; (b) max/min on $x^2+y^2\le4$.
\end{exercise}
\begin{solbox}
\textbf{(a)} $T_x=2x-y=0$, $T_y=-x+2y=0\Rightarrow x=y=0$.
$D=T_{xx}T_{yy}-T_{xy}^2=2\cdot2-(-1)^2=3>0$, $T_{xx}>0$: global \textbf{minimum} at $(0,0)$, $T=0$.

\textbf{(b)} Interior: $T(0,0)=0$.
Boundary $x=2\cos\theta$, $y=2\sin\theta$:
\[
  T=4\cos^2\!\theta-4\cos\theta\sin\theta+4\sin^2\!\theta=4-2\sin(2\theta).
\]
Max when $\sin(2\theta)=-1$: $T=6$ at $\theta=3\pi/4$, i.e.\ $(-\sqrt{2},\sqrt{2})$ (and $(\sqrt{2},-\sqrt{2})$).
Min when $\sin(2\theta)=1$: $T=2$ at $\theta=\pi/4$, i.e.\ $(\sqrt{2},\sqrt{2})$.

Overall: \textbf{Absolute max} $=6$, \textbf{absolute min} $=0$.
\end{solbox}

%------------------------------------------------------------------
\begin{exercise}
$f(x,y)=x^2+4xy+y^3$. (a) $dy/dx$ if $f=0$; (b) $\grad f$ at $(1,-1)$;
(c) directional derivative at $(1,-1)$ toward origin.
\end{exercise}
\begin{solbox}
\textbf{(a)} By implicit differentiation:
\[
  \frac{dy}{dx}=-\frac{F_x}{F_y}=-\frac{2x+4y}{4x+3y^2}.
\]

\textbf{(b)} $\grad f=\langle 2x+4y,\;4x+3y^2\rangle$.
At $(1,-1)$: $\grad f=\langle2-4,\;4+3\rangle=\langle-2,7\rangle$.

\textbf{(c)} Direction toward origin from $(1,-1)$: $\mathbf{d}=\langle-1,1\rangle$, $|\mathbf{d}|=\sqrt{2}$.
Unit vector $\mathbf{u}=\langle-1,1\rangle/\sqrt{2}$.
\[
  D_{\mathbf{u}}f=\langle-2,7\rangle\cdot\frac{\langle-1,1\rangle}{\sqrt{2}}
  =\frac{2+7}{\sqrt{2}}=\frac{9\sqrt{2}}{2}.
\]
\end{solbox}

%------------------------------------------------------------------
\begin{exercise}
Evaluate $I=\int_{-\infty}^{\infty}e^{-x^2}\dif{x}$ via polar coordinates.
\end{exercise}
\begin{solbox}
\[
  I^2=\iint_{\R^2}e^{-(x^2+y^2)}\dif{A}
  =\int_0^{2\pi}\!\int_0^\infty e^{-r^2}r\dif{r}\dif{\theta}
  =2\pi\left[-\frac{e^{-r^2}}{2}\right]_0^\infty=2\pi\cdot\frac{1}{2}=\pi.
\]
Therefore $I=\sqrt{\pi}$.\hfill$\square$
\end{solbox}

%------------------------------------------------------------------
\begin{exercise}
Volume between paraboloid $z=x^2+y^2$ and sphere $x^2+y^2+z^2=6$.
\end{exercise}
\begin{solbox}
\emph{Choice of coordinates:} cylindrical, since the paraboloid is $z=r^2$ and sphere is $r^2+z^2=6$.
Intersection: $r^2+(r^2)^2=6\Rightarrow r^4+r^2-6=0\Rightarrow(r^2-2)(r^2+3)=0\Rightarrow r=\sqrt{2}$, $z=2$.
\[
  V=\int_0^{2\pi}\!\int_0^{\sqrt{2}}\!\int_{r^2}^{\sqrt{6-r^2}}r\dif{z}\dif{r}\dif{\theta}
  =2\pi\int_0^{\sqrt{2}}r\!\left(\sqrt{6-r^2}-r^2\right)\dif{r}.
\]
For $\int_0^{\sqrt{2}}r\sqrt{6-r^2}\dif{r}$: let $u=6-r^2$:
$=\left[-\tfrac{(6-r^2)^{3/2}}{3}\right]_0^{\sqrt{2}}=-\tfrac{4^{3/2}}{3}+\tfrac{6^{3/2}}{3}=\tfrac{6\sqrt{6}-8}{3}$.
For $\int_0^{\sqrt{2}}r^3\dif{r}=\left[\tfrac{r^4}{4}\right]_0^{\sqrt{2}}=1$.
\[
  V=2\pi\!\left(\frac{6\sqrt{6}-8}{3}-1\right)=\frac{2\pi(6\sqrt{6}-11)}{3}.
\]
\end{solbox}

%------------------------------------------------------------------
\begin{exercise}
Curve $F(x,y)=x\sin(y)-y\cos(x)=0$.
(a) Find $dy/dx$ near $(0,0)$; (b) tangent line at $(0,0)$.
\end{exercise}
\begin{solbox}
$F_x=\sin(y)+y\sin(x)$, $F_y=x\cos(y)-\cos(x)$.
At $(0,0)$: $F_x=0$, $F_y=0-1=-1$.
\[
  \frac{dy}{dx}=-\frac{F_x}{F_y}=-\frac{0}{-1}=0.
\]
\textbf{(b)} Slope $=0$ at $(0,0)$, and $F(0,0)=0\cdot\sin0-0\cdot\cos0=0$ (point is on the curve).
Tangent line: $\boxed{y=0}$.
\end{solbox}

\bigskip
\begin{center}
  \rule{0.6\linewidth}{0.4pt}\\[4pt]
  {\small\itshape End of Solutions}
\end{center}

\end{document}
